\begin{helper}
Fix a projection-size history cell \(C_h\) with recorded coordinate set
\(R_h\) and terminal coordinate set \(U_h\).  Let \(\omega^\ast\in C_h\) be
a completed sample used only to encode the already exposed first-color
realization.  The same-color terminal values of \(\omega^\ast\) are ignored.
Assume
\[
0\le\varepsilon_1,\qquad
\noderef{ParameterHierarchyEventualComparisons}
(\varepsilon,\varepsilon_1,\varepsilon_2,n_0),\qquad n>n_0,
\]
\[
m=\noderef{TwoBiteNaturalReducedVertexCount}(n),\qquad
p=\noderef{TwoBiteEdgeProbability}\!\left(\frac12,n\right),
\]
and
\[
|I|=k=\noderef{TwoBiteNaturalIndependenceScale}(1+\varepsilon,n).
\]
Assume also that the history package identifies the recorded and terminal
coordinates:
\[
e\in R_h
\quad\Longleftrightarrow\quad
e\in\noderef{TwoBitePreTerminalRecordedPairs}(\omega,\varepsilon,I)
\qquad(\omega\in C_h),
\tag{H1}
\]
and
\[
e\in U_h
\quad\Longleftrightarrow\quad
e\in\noderef{TwoBiteTerminalCoordinateUniverse}(m)
\hbox{ and }e\notin R_h.
\tag{H2}
\]
Finally, assume that every staged-open coordinate in every completion of the
history belongs to \(U_h\).

There are two finite candidate sets \(E_R^\ast,E_B^\ast\subseteq U_h\).  The
red candidate set is defined by the exposed blue realization:
\[
e\in E_R^\ast
\]
if and only if \(e=\operatorname{red}(q)\in U_h\) and there is a blue base
vertex \(b\) which is large, medium, or small for \(I\) in \(\omega^\ast\)
such that
\[
q\in\noderef{TwoBiteBasePairs}
\bigl(\pi_R^{\omega^\ast}
(\noderef{TwoBiteX}(\omega^\ast,I,\operatorname{blue}(b)))\bigr).
\]
The blue candidate set is defined symmetrically from the exposed red
realization:
\[
e\in E_B^\ast
\]
if and only if \(e=\operatorname{blue}(q)\in U_h\) and there is a red base
vertex \(r\) which is large, medium, or small for \(I\) in \(\omega^\ast\)
such that
\[
q\in\noderef{TwoBiteBasePairs}
\bigl(\pi_B^{\omega^\ast}
(\noderef{TwoBiteX}(\omega^\ast,I,\operatorname{red}(r)))\bigr).
\]

If a completion \(\omega\in C_h\) has the same injection and the same
recorded values as \(\omega^\ast\), agrees with \(\omega^\ast\) on all blue
terminal coordinates, and satisfies
\[
\noderef{TwoBiteRegularityEvent}(k,\varepsilon,\varepsilon_1,\varepsilon_2,
\omega),
\]
then
\[
\noderef{ClosedPairsBound}(|E_R^\ast|,3\varepsilon_1,k).
\]
Moreover, if \(I\) is independent in the shadow of
\(\noderef{TwoBiteFinalGraph}(\omega)\), then every present red
staged-open coordinate of \(\omega\) belongs to \(E_R^\ast\).  The symmetric
statement holds with the colors interchanged: if \(\omega\) agrees with
\(\omega^\ast\) on the red terminal coordinates and is regular, then
\[
\noderef{ClosedPairsBound}(|E_B^\ast|,3\varepsilon_1,k),
\]
and every present blue staged-open coordinate required to be deleted for the
independence of \(I\) belongs to \(E_B^\ast\).
\end{helper}
\begin{proof}
We define \(E_R^\ast\) and \(E_B^\ast\) by the two displayed membership
rules.  They are finite because both are filters of the finite terminal set
\(U_h\).  The red rule uses only the injection of \(\omega^\ast\), the blue
edge values of \(\omega^\ast\), and the terminal set \(U_h\): the set
\(\noderef{TwoBiteX}(\omega^\ast,I,\operatorname{blue}(b))\) asks only
which vertices of \(I\) are blue-neighbors of \(b\), and the projection
\(\pi_R^{\omega^\ast}\) asks only for the fixed injection.  No red terminal
edge value is queried.  This last assertion is exactly the blue side of
\(\noderef{FixedSetOppositeColorTerminalCoverage}\): any blue base-edge
coordinate needed for such a neighborhood is either in \(R_h\), where the
history has already fixed it, or in the full unrecorded terminal set \(U_h\),
where it is one of the blue first-color values.  Thus once the blue terminal
realization has been exposed, \(E_R^\ast\) is fixed before any red terminal
Bernoulli variables are sampled.  The blue statement is identical after
swapping the two colors.

We prove the red bound and containment; the blue proof is symmetric.  Let
\(\omega\in C_h\) be a completion with the same injection and recorded table
as \(\omega^\ast\), with the same blue values on \(U_h\), and satisfying the
regularity event.  Apply
\(\noderef{FixedSetOppositeColorTerminalCoverage}\) with the data
\[
I,\ R_h,\ U_h,\ \omega,\ \omega^\ast .
\]
Its hypotheses are all present.  The common-injection assumption is one of
the hypotheses of the present continuation.  Agreement on \(R_h\) is the
recorded-table hypothesis.  The exact pre-terminal recorded-set hypotheses
for both \(\omega\) and \(\omega^\ast\) follow from (H1), since both samples
lie in \(C_h\).  The full-terminal-universe hypothesis is (H2).  Finally,
the first-color terminal hypothesis says precisely that the two samples agree
on every blue coordinate in \(U_h\).  The blue side of the child theorem
therefore gives, for every blue base vertex \(b\),
\[
\noderef{TwoBiteX}(\omega,I,\operatorname{blue}(b))
=
\noderef{TwoBiteX}(\omega^\ast,I,\operatorname{blue}(b)), \tag{1}
\]
the equivalence of the large, medium, and small classifications of
\(\operatorname{blue}(b)\) in the two samples, and the projected-pair
equivalence
\[
q\in\noderef{TwoBiteBasePairs}
\bigl(\pi_R^\omega
(\noderef{TwoBiteX}(\omega,I,\operatorname{blue}(b)))\bigr)
\Longleftrightarrow
q\in\noderef{TwoBiteBasePairs}
\bigl(\pi_R^{\omega^\ast}
(\noderef{TwoBiteX}(\omega^\ast,I,\operatorname{blue}(b)))\bigr).
\tag{2}
\]
Consequently the displayed red-exceptional predicate defining \(E_R^\ast\)
is exactly the red-exceptional predicate one would define from \(\omega\),
except that it is restricted to \(U_h\).

It remains to count this predicate.  Split \(E_R^\ast\) according to whether
the chosen blue witness is large, medium, or small.  If a coordinate
\(\operatorname{red}(q)\in E_R^\ast\) is certified by a witness in one of
these three classes, choose one such witness \(b\) and one class label.  The
defining membership of \(E_R^\ast\) gives the projected-pair condition in
\(\omega^\ast\).  By (2) the same \(q\) satisfies the corresponding
projected-pair condition in \(\omega\).  By the classification equivalences
from the coverage child, the same \(b\) lies in the same large, medium, or
small class for \(\omega\).  Therefore, for that class, the final pair of
vertices of \(I\) producing \(q\) belongs to
\[
\noderef{TwoBiteClosedPairsInClass}(\omega,I,\mathcal A).
\tag{3}
\]
The map from red tagged coordinates to final unordered pairs has
multiplicity at most one, because the red projection of a fixed unordered
final pair determines a unique increasing base pair in
\(\noderef{TwoBiteBasePairs}\).  Hence the number of red candidate
coordinates certified by class \(\mathcal A\) is at most the number of
closed pairs in that class.

For the large class, \(\noderef{LargeClosedPairsBound}\), together with the
eventual hierarchy, \(n>n_0\), the rounded parameter identities, the scale
identity for \(I\), and the fiber-degree part of regularity, gives a
\(\noderef{ClosedPairsBound}(-,\varepsilon_1,k)\) bound.  For the medium and
small classes, the medium and small closed-pair conjuncts of
\(\noderef{TwoBiteRegularityEvent}\) give the same
\(\varepsilon_1k^2\) bound.  Summing the three class contributions gives
\[
|E_R^\ast|\le
\varepsilon_1k^2+\varepsilon_1k^2+\varepsilon_1k^2
=3\varepsilon_1k^2,
\]
which is
\(\noderef{ClosedPairsBound}(|E_R^\ast|,3\varepsilon_1,k)\).
This proves the promised size bound uniformly for every regular completion
with the same exposed blue realization.  Notice that the bound is for the
fixed set \(E_R^\ast\), not for a set chosen after red terminal values are
known.

Now assume in addition that \(I\) is independent in the final shadow graph of
\(\omega\), and let \(e=\operatorname{red}(q)\) be a present red
staged-open coordinate of \(\omega\).  By the history hypothesis on staged
open coordinates, \(e\in U_h\).  Apply
\(\noderef{FixedSetExceptionalWitnessClassification}\) to the completed
regular sample \(\omega\).  Its hypotheses are exactly the scalar
hypotheses above, the regularity of \(\omega\), and the independence of
\(I\).  The child theorem says that this present staged-open red coordinate
has a non-huge opposite-color blue witness: there is a blue vertex \(b\) that
is large, medium, or small for \(I\) in \(\omega\), and
\[
q\in\noderef{TwoBiteBasePairs}
\bigl(\pi_R^\omega
(\noderef{TwoBiteX}(\omega,I,\operatorname{blue}(b)))\bigr).
\tag{4}
\]
Using (2), the projected-pair condition in (4) transfers to
\(\omega^\ast\); using the classification equivalences from the coverage
child, the same \(b\) is large, medium, or small for
\(\omega^\ast\).  Thus \(b\) certifies the displayed membership condition
for \(E_R^\ast\).  Since \(e\in U_h\), the defining equivalence of
\(E_R^\ast\) gives \(e\in E_R^\ast\).

The blue case repeats the argument with the colors interchanged.  Agreement
on the red terminal coordinates fixes all red neighborhoods and red
large/medium/small classifications before blue terminal values are exposed;
this is the red side of
\(\noderef{FixedSetOppositeColorTerminalCoverage}\), using (H1), (H2), the
common injection, and recorded agreement in exactly the same way.  The same
three class bounds give
\(\noderef{ClosedPairsBound}(|E_B^\ast|,3\varepsilon_1,k)\); and
\(\noderef{FixedSetExceptionalWitnessClassification}\), transferred through
the common injection and red first-color agreement, puts every present blue
staged-open coordinate needed for final independence into \(E_B^\ast\).
\end{proof}
